\documentclass[a4paper,11pt]{article}

% ============================================================
%  Le Collé — Séance 1 : Déterminants, courbes paramétrées
%  © 2026 Ahmed Amine Ouqlani
% ============================================================

% ---------- Polices ----------
\usepackage[utf8]{inputenc}
\usepackage[T1]{fontenc}
\usepackage[sc]{mathpazo}          % texte et mathématiques : Palatino
\usepackage[scaled=0.88]{helvet}   % titres et repères : linéale
\linespread{1.07}
\usepackage{amsmath,amssymb,amsthm,mathrsfs,stmaryrd}
\usepackage{microtype}

% ---------- Mise en page ----------
\usepackage[a4paper,top=2.3cm,bottom=2.3cm,left=2.4cm,right=2.4cm,
            headheight=14pt,headsep=0.8cm,footskip=1.1cm]{geometry}
\setlength{\parindent}{0pt}
\setlength{\parskip}{0.45em plus 0.15em minus 0.1em}
\raggedbottom
\emergencystretch=1.5em
\clubpenalty=10000 \widowpenalty=10000 \displaywidowpenalty=10000 \brokenpenalty=10000

% ---------- Couleurs : un bleu de nuit, un ocre, des gris ----------
\usepackage{xcolor}
\definecolor{Ink}{HTML}{1C2533}
\definecolor{Bleu}{HTML}{1F3A5F}
\definecolor{BleuMoyen}{HTML}{4A6A8F}
\definecolor{BleuClair}{HTML}{9DB2C9}
\definecolor{Fond}{HTML}{F2F4F7}
\definecolor{Ocre}{HTML}{B07A26}
\definecolor{Brique}{HTML}{A23B2A}
\definecolor{Gris}{HTML}{5E6672}
\definecolor{Filet}{HTML}{C5CCD6}

% ---------- Packages ----------
\usepackage{enumitem}
\setlist{itemsep=0.25em,topsep=0.35em,parsep=0pt}
\setlist[itemize]{label=\textcolor{BleuMoyen}{\small\textbullet},leftmargin=1.4em}
\setlist[enumerate]{leftmargin=1.8em,label=\textsf{\small\color{Bleu}\arabic*.}}
\usepackage{booktabs,array,multicol}
\usepackage{needspace,float}
\usepackage[font=small,labelfont={sf,color=Bleu},labelsep=period,skip=6pt]{caption}
\usepackage{tikz}
\usetikzlibrary{arrows.meta,calc,decorations.markings}
\usepackage{pgfplots}
\pgfplotsset{compat=1.17}
\usepackage{tkz-tab}
\usepackage[most]{tcolorbox}
\usepackage{titlesec,titletoc}
\usepackage{fancyhdr,lastpage,etoolbox}
\usepackage[colorlinks=true,linkcolor=Bleu,urlcolor=Bleu,citecolor=Bleu,
            pdftitle={Le Collé — Équations différentielles : sélection de 5 problèmes},pdfauthor={Ahmed Amine Ouqlani}]{hyperref}

\numberwithin{equation}{section}
\numberwithin{figure}{section}
\numberwithin{table}{section}

% ---------- Titres ----------
% Chapitres (sections) : grand numéro, titre en linéale, filet ocre
\titleformat{\section}[display]
  {\normalfont\sffamily\color{Ink}}
  {{\small\color{Gris}\textls[200]{\MakeUppercase{Chapitre}}}\\[-0.2em]%
   {\fontsize{64}{64}\selectfont\color{BleuClair}\thesection}}
  {0.4em}
  {\fontsize{25}{30}\selectfont\bfseries}
  [\vspace{0.5em}{\color{Ocre}\rule{3.2cm}{1.4pt}}]
\titlespacing*{\section}{0pt}{0pt}{1.6em}
\titleformat{name=\section,numberless}[display]
  {\normalfont\sffamily\color{Ink}}{}{0pt}{\fontsize{25}{30}\selectfont\bfseries}
  [\vspace{0.5em}{\color{Ocre}\rule{3.2cm}{1.4pt}}]
\titlespacing*{name=\section,numberless}{0pt}{0pt}{1.6em}

\titleformat{\subsection}
  {\needspace{5\baselineskip}\normalfont\sffamily\Large\bfseries\color{Ink}}
  {\color{Bleu}\thesubsection}{0.8em}{}
\titlespacing*{\subsection}{0pt}{1.8em}{0.5em}
\titleformat{\subsubsection}
  {\needspace{5\baselineskip}\normalfont\sffamily\large\bfseries\color{Ink}}
  {\color{Bleu}\thesubsubsection}{0.7em}{}
\titlespacing*{\subsubsection}{0pt}{1.6em}{0.4em}
\titleformat{\paragraph}[runin]{\normalfont\sffamily\bfseries\color{Ink}}{}{0pt}{}[.]
\titlespacing*{\paragraph}{0pt}{0.9em}{0.6em}
\setcounter{secnumdepth}{3}
\setcounter{tocdepth}{2}

% ---------- Table des matières ----------
\titlecontents{section}[0pt]{\addvspace{0.7em}\sffamily\bfseries\color{Ink}}
  {\makebox[2.2em][l]{\color{Bleu}\thecontentslabel}}{}{\hfill\contentspage}[\addvspace{0.2em}]
\titlecontents{subsection}[2.2em]{\small\linespread{1}\selectfont}
  {\makebox[2.6em][l]{\color{Gris}\thecontentslabel}}{}{\titlerule*[0.7pc]{.}\contentspage}

% ---------- En-têtes et pieds de page ----------
\pagestyle{fancy}
\fancyhf{}
\renewcommand{\sectionmark}[1]{\markboth{\thesection\quad#1}{}}
\fancyhead[L]{\footnotesize\sffamily\color{Gris}\leftmark}
\fancyhead[R]{\footnotesize\sffamily\scshape\color{Gris}Le Collé}
\fancyfoot[C]{\small\sffamily\color{Ink}\thepage}
\fancyfoot[L]{\scriptsize\sffamily\hyperlink{toc}{\color{Filet}sommaire}}
\renewcommand{\headrulewidth}{0.4pt}
\renewcommand{\headrule}{{\color{Filet}\hrule height\headrulewidth width\headwidth}}
\renewcommand{\footrulewidth}{0pt}
\fancypagestyle{plain}{\fancyhf{}\fancyfoot[C]{\small\sffamily\color{Ink}\thepage}\renewcommand{\headrulewidth}{0pt}}

% ---------- Énoncés ----------
\newtheoremstyle{enonce}{0pt}{0pt}{}{}{\sffamily\bfseries\color{Bleu}}{.}{0.5em}
  {\thmname{#1}\thmnumber{~#2}\thmnote{\ \normalfont\sffamily\color{Ink}(#3)}}
\newtheoremstyle{exemple}{0.9em}{0.9em}{}{}{\sffamily\bfseries\color{BleuMoyen}}{.}{0.5em}
  {\thmname{#1}\thmnumber{~#2}\thmnote{\ \normalfont\sffamily\color{Ink}(#3)}}
\newtheoremstyle{remarque}{0.7em}{0.7em}{}{}{\sffamily\itshape}{.}{0.5em}{\thmname{#1}\thmnote{\ (#3)}}
\theoremstyle{enonce}
\newtheorem{definition}{Définition}[section]
\newtheorem{theoreme}[definition]{Théorème}
\newtheorem{proposition}[definition]{Proposition}
\newtheorem{corollaire}[definition]{Corollaire}
\theoremstyle{exemple}
\newtheorem{exemple}[definition]{Exemple}
\theoremstyle{remarque}
\newtheorem*{remarque}{Remarque}

\tcbset{
  defi/.style={enhanced,breakable,colback=white,colframe=white,boxrule=0pt,arc=0pt,
     borderline west={1.2pt}{0pt}{Bleu},left=11pt,right=4pt,top=3pt,bottom=3pt,
     before skip=1.1em,after skip=1.1em},
  thm/.style={enhanced,breakable,colback=Fond,colframe=Fond,boxrule=0pt,arc=0pt,
     borderline west={2.4pt}{0pt}{Bleu},left=11pt,right=8pt,top=6pt,bottom=6pt,
     before skip=1.2em,after skip=1.2em}}
\tcolorboxenvironment{definition}{defi}
\tcolorboxenvironment{theoreme}{thm}
\tcolorboxenvironment{proposition}{thm}
\tcolorboxenvironment{corollaire}{thm}

\renewcommand{\proofname}{Démonstration}
\renewcommand{\qedsymbol}{\textcolor{BleuClair}{$\blacksquare$}}
\makeatletter
\renewenvironment{proof}[1][\proofname]{\par\pushQED{\qed}\normalfont
  \topsep0.5em\relax\trivlist\item[\hskip\labelsep\sffamily\itshape\color{Gris}#1.]\ignorespaces}
  {\popQED\endtrivlist\@endpefalse}
\makeatother

% ---------- Méthode : filets haut et bas, titre en bandeau ----------
\makeatletter\newcommand{\methlabelfix}{\protected@edef\@currentlabel{\thetcbcounter}}\makeatother
\newtcolorbox[auto counter,number within=section]{methode}[2][]{%
  enhanced,unbreakable,colback=white,colframe=white,boxrule=0pt,arc=0pt,
  borderline north={0.9pt}{0pt}{Bleu},borderline south={0.9pt}{0pt}{Bleu},
  left=4pt,right=4pt,top=8pt,bottom=8pt,before skip=1.4em,after skip=1.4em,
  before upper={\methlabelfix\phantomsection\addcontentsline{met}{metentry}{\protect\makebox[2.6em][l]{\thetcbcounter}#2}%
     {\sffamily\bfseries\small\color{white}\colorbox{Bleu}{\strut\ MÉTHODE~\thetcbcounter\ }}%
     \ \ {\sffamily\bfseries\color{Bleu}#2}\par\smallskip},#1}

% ---------- Repères pédagogiques : un filet de couleur, un mot-clé ----------
\newcommand{\definirbloc}[3]{% nom, couleur, libellé
  \newtcolorbox{#1}[1][]{enhanced,breakable,blanker,
    borderline west={1.4pt}{0pt}{#2},left=10pt,right=2pt,top=1pt,bottom=1pt,
    before skip=1em,after skip=1em,
    before upper={{\sffamily\bfseries\footnotesize\color{#2}\textls[120]{\MakeUppercase{#3}}}\quad},##1}}
\definirbloc{reflexe}{Bleu}{Réflexe}
\definirbloc{astuce}{Ocre}{Astuce}
\definirbloc{attention}{Brique}{Attention}
\definirbloc{enpratique}{BleuMoyen}{En pratique}
\definirbloc{plusloin}{Gris}{Pour aller plus loin}
\definirbloc{strategie}{Bleu}{Stratégie}
\definirbloc{lecture}{BleuMoyen}{Lire le schéma}
% compatibilité avec les parties reprises
\newenvironment{garde}{\begin{attention}}{\end{attention}}
\newenvironment{complementHP}{\begin{plusloin}}{\end{plusloin}}
\newenvironment{interpretation}{\begin{enpratique}}{\end{enpratique}}

% ---------- Exercices, corrigés ----------
\newcounter{exo}[section]
\renewcommand{\theexo}{\thesection.\arabic{exo}}
\newcommand{\niveau}[1]{{\color{Ocre}\footnotesize
  \ifcase#1\or$\bullet\circ\circ\circ\circ$\or$\bullet\bullet\circ\circ\circ$\or$\bullet\bullet\bullet\circ\circ$%
  \or$\bullet\bullet\bullet\bullet\circ$\or$\bullet\bullet\bullet\bullet\bullet$\fi}}
\newcommand{\serie}[1]{\par\needspace{8\baselineskip}\vspace{1.5em}%
  \noindent{\sffamily\bfseries\color{Bleu}#1}\par\nobreak\vspace{-0.4em}%
  \noindent{\color{Filet}\rule{\linewidth}{0.4pt}}\par\nobreak\vspace{0.1em}}
\newenvironment{exo}[3]{%
  \par\needspace{5\baselineskip}\vspace{0.75em}\refstepcounter{exo}\label{e:#1}%
  \def\exolab{#1}%
  \noindent{\sffamily\bfseries\color{Bleu}Exercice~\theexo}\quad\niveau{#2}%
  \hfill{\footnotesize\sffamily\color{Gris}#3}\par\nobreak\vspace{0.15em}}%
  {\par\nopagebreak\noindent\hfill{\scriptsize\sffamily\hyperref[c:\exolab]{\color{Gris}corrigé p.~\pageref{c:\exolab}}}\par}
\newenvironment{corrige}[1]{%
  \par\needspace{6\baselineskip}\vspace{1em}\phantomsection\label{c:#1}%
  \noindent{\sffamily\bfseries\color{Bleu}Exercice~\ref{e:#1}}\quad
  {\scriptsize\sffamily\hyperref[e:#1]{\color{Gris}énoncé p.~\pageref{e:#1}}}\par\nobreak\vspace{0.2em}}%
  {\hfill\textcolor{BleuClair}{$\square$}\par}
\newcommand{\etape}[1]{\par\smallskip\noindent{\sffamily\bfseries\small\color{BleuMoyen}#1.}\ }
\newenvironment{solution}{\par\noindent{\sffamily\bfseries\small\color{BleuMoyen}Solution.}\ }{\hfill\textcolor{BleuClair}{$\square$}\par}
\newcommand{\chercher}{\par\smallskip\noindent{\color{Filet}\rule[0.5ex]{1.2cm}{0.4pt}}\
  {\footnotesize\sffamily\itshape\color{Gris}laisser chercher l'étudiant avant de poursuivre}\par\smallskip}

% ---------- Guide de séance ----------

\newcommand{\marge}[1]{\leavevmode\marginpar{\raggedleft\scriptsize\sffamily\color{Gris}\textls[100]{\MakeUppercase{#1}}}\ignorespaces}
\newcommand{\Etape}[4]{%
  \subsection[Étape #1 --- #3 {\color{Gris}(#2)}]{#3}\label{etape:#1}%
  \noindent{\small\sffamily\color{Gris}{\color{Ocre}\bfseries Étape #1}\quad·\quad #2\quad·\quad #4}\par\medskip}
\newcommand{\question}[1]{\par\smallskip\noindent{\sffamily\bfseries\small\color{BleuMoyen}Question.}\ \textit{#1}\par\smallskip}
\newcommand{\rappel}[1]{\par\smallskip\noindent\phantomsection\addcontentsline{rap}{rapentry}{#1}%
  {\sffamily\bfseries\scriptsize\colorbox{Fond}{\color{Gris}\strut RAPPEL}}\ {\sffamily\small\color{Gris}#1.}\ }
\newcommand{\pourquoi}{\par\noindent{\sffamily\bfseries\small\color{Ocre}Pourquoi ?}\ }

% Ouverture de chapitre : \ouverture{objectifs}{prérequis}{durée}
\newcommand{\ouverture}[3]{%
  \begin{tcolorbox}[enhanced,colback=Fond,colframe=Fond,arc=0pt,boxrule=0pt,left=12pt,right=12pt,top=10pt,bottom=10pt,
     before skip=1.2em,after skip=1.6em]
  \begin{minipage}[t]{0.56\linewidth}
  {\sffamily\bfseries\small\color{Bleu}OBJECTIFS}\par\vspace{0.3em}\small #1
  \end{minipage}\hfill
  \begin{minipage}[t]{0.38\linewidth}
  {\sffamily\bfseries\small\color{Bleu}PRÉREQUIS}\par\vspace{0.3em}\small #2
  \ifx\relax#3\relax\else\par\vspace{0.6em}{\sffamily\bfseries\small\color{Bleu}DURÉE}\par\vspace{0.2em}\small #3\fi
  \end{minipage}
  \end{tcolorbox}}
% Fin de chapitre
\newenvironment{essentiel}{\par\needspace{12\baselineskip}\vspace{1.5em}%
  \begin{tcolorbox}[enhanced,unbreakable,colback=white,colframe=white,boxrule=0pt,arc=0pt,
   borderline north={1.4pt}{0pt}{Ocre},left=2pt,right=2pt,top=8pt,bottom=4pt]
  {\sffamily\bfseries\color{Ocre}L'ESSENTIEL DU CHAPITRE}\par\medskip}{\end{tcolorbox}}

% ---------- Notations ----------
\newcommand{\K}{\mathbb{K}}
\newcommand{\R}{\mathbb{R}}
\newcommand{\C}{\mathbb{C}}
\newcommand{\N}{\mathbb{N}}
\newcommand{\Z}{\mathbb{Z}}
\newcommand{\Mn}[1]{\mathcal{M}_{n}(#1)}
\newcommand{\GLn}[1]{\mathrm{GL}_n(#1)}
\newcommand{\LE}{\mathcal{L}(E)}
\newcommand{\rg}{\operatorname{rg}}
\newcommand{\vect}{\operatorname{Vect}}
\newcommand{\Mat}{\operatorname{Mat}}
\newcommand{\Id}{\operatorname{Id}}
\newcommand{\sh}{\operatorname{sh}}
\newcommand{\ch}{\operatorname{ch}}
\renewcommand{\th}{\operatorname{th}}
\DeclareMathOperator{\argsh}{argsh}
\newcommand{\vv}[1]{\overrightarrow{#1}}
\newcommand{\ve}[1]{\vec{#1}}
\newcommand{\ient}[2]{\llbracket #1,#2\rrbracket}
\newcommand{\dt}[1]{\left|\begin{matrix}#1\end{matrix}\right|}
\newcommand{\dis}{\displaystyle}

% ---------- Figures ----------
\pgfplotsset{trig format plots=rad}
\pgfplotsset{courbe/.style={axis lines=middle,axis equal image,
  axis line style={Gris,-{Stealth[length=5pt]}},
  xlabel={$x$},ylabel={$y$},ticklabel style={font=\scriptsize\sffamily,color=Gris},
  every axis x label/.style={at={(ticklabel* cs:1)},anchor=west,font=\small},
  every axis y label/.style={at={(ticklabel* cs:1)},anchor=south,font=\small},
  samples=300,smooth,width=0.6\linewidth}}
\tikzset{arcetude/.style={Bleu,line width=1.25pt},
         reste/.style={BleuClair,line width=1pt},
         tgte/.style={Ocre,line width=0.9pt},
         asy/.style={Gris,dashed,line width=0.6pt},
         pt/.style={circle,fill=Ink,inner sep=1.2pt}}

% ---------- Navigation ----------
\makeatletter
\newcommand*\l@rapentry{\@dottedtocline{1}{0em}{0em}}
\newcommand*\l@metentry{\@dottedtocline{1}{0em}{0em}}
\newcommand{\listemethodes}{\@starttoc{met}}
\newcommand{\listerappels}{\@starttoc{rap}}
\makeatother
% onglets de chapitre sur le bord droit
\usepackage{eso-pic}
\newif\ifonglets \ongletsfalse
\newcommand{\dessineronglet}{%
  \edef\ongcol{\ifnum\value{section}=3 Ocre\else Bleu\fi}%
  \edef\ongnum{\the\value{section}}%
  \pgfmathsetmacro{\ongy}{3.2+2.3*(\ongnum-1)}%
  \begin{tikzpicture}[remember picture,overlay]
  \fill[\ongcol] ([xshift=-0.75cm,yshift=-\ongy cm]current page.north east)
     rectangle ([yshift=-\ongy cm-1.9cm]current page.north east);
  \node[text=white,font=\sffamily\bfseries\large] at ([xshift=-0.37cm,yshift=-\ongy cm-0.95cm]current page.north east)
     {\ongnum};
  \end{tikzpicture}}
\AddToShipoutPictureBG{\ifonglets\ifnum\value{section}>0 \dessineronglet\fi\fi}
% blocs de navigation et d'enseignement
\definirbloc{commencer}{Gris}{Si vous commencez ici}
\definirbloc{horsprogramme}{Brique}{Hors programme}
\newcommand{\HP}{\,{\sffamily\bfseries\scriptsize\textcolor{Brique}{[HP]}}}
\newcommand{\source}[1]{{\footnotesize\sffamily\color{Gris}\,(#1)}}
\newcommand{\coupdepouce}[1]{\par\vspace{-0.2em}\noindent{\footnotesize\sffamily\color{Gris}%
  \textbf{Coup de pouce} (à donner si l'étudiant bloque) : #1}\par}
\newcommand{\pointarret}[1]{\par\vspace{1em}\noindent{\color{Filet}\rule[0.5ex]{\linewidth}{0.4pt}}\par\nobreak
  \noindent{\footnotesize\sffamily\color{Ocre}\bfseries POINT D'ARRÊT POSSIBLE}\quad{\footnotesize\sffamily\color{Gris}#1}\par\vspace{0.6em}}
\newcommand{\demarche}[1]{\noindent{\sffamily\bfseries\small\color{BleuMoyen}Démarche.}\ \textit{#1}\par\smallskip}

% ---------- Corrigés : idées, astuces, pièges, à retenir ----------
\newcommand{\choix}[1]{\par\noindent{\sffamily\bfseries\small\color{BleuMoyen}Choix de la méthode.}\ #1\par\smallskip}
\newcommand{\gain}[1]{\par\noindent{\sffamily\bfseries\small\color{Ocre}Astuce.}\ #1\par}
\newcommand{\piege}[1]{\par\noindent{\sffamily\bfseries\small\color{Brique}Piège.}\ #1\par}
\newtcolorbox{retenirbox}{enhanced,breakable,colback=Fond,colframe=Fond,arc=0pt,boxrule=0pt,
  borderline west={2pt}{0pt}{Ocre},left=8pt,right=6pt,top=4pt,bottom=4pt,before skip=0.7em,after skip=0.5em}
\newcommand{\retenir}[1]{\begin{retenirbox}{\sffamily\bfseries\footnotesize\color{Ocre}À RETENIR}\quad #1\end{retenirbox}}
\newcommand{\TD}[1]{{\color{Ocre}\bfseries TD\,#1}\ ·\ }
\newcommand{\q}[1]{\par\smallskip\noindent{\sffamily\bfseries\color{Bleu}#1}\ }
\providecommand{\iddots}{\mathinner{\mkern1mu\raise1pt\hbox{.}\mkern2mu\raise4pt\hbox{.}\mkern2mu\raise7pt\hbox{.}\mkern1mu}}

\providecommand{\Ker}{\operatorname{Ker}}\renewcommand{\Im}{\operatorname{Im}}\providecommand{\tr}{\operatorname{tr}}\providecommand{\Sp}{\operatorname{Sp}}

% ================================================================ TD : mise en page propre au livret d'exercices
\fancyhf{}
\fancyhead[L]{\footnotesize\sffamily\color{Gris}Sélection · Équations différentielles}
\fancyhead[R]{\footnotesize\sffamily\scshape\color{Gris}Le Collé}
\fancyfoot[C]{\small\sffamily\color{Ink}\thepage\,/\,\pageref{LastPage}}
\fancypagestyle{plain}{\fancyhf{}\fancyfoot[C]{\small\sffamily\color{Ink}\thepage\,/\,\pageref{LastPage}}\renewcommand{\headrulewidth}{0pt}}
\makeatletter
\newcommand{\partie}[3]{% numéro, titre, description
  \par\needspace{10\baselineskip}\vspace{1.4em}\phantomsection\label{partie:#1}\pdfbookmark[0]{Partie #1}{partie.#1}%
  {\sffamily\small\color{Gris}\textls[200]{PARTIE #1}}\par\nobreak\vspace{0.1em}
  {\sffamily\bfseries\LARGE\color{Ink}#2}\par\nobreak\vspace{0.3em}
  {\color{Ocre}\rule{2.6cm}{1.2pt}}\par\nobreak\vspace{0.2em}
  {\small\color{Gris}#3}\par\nobreak\vspace{0.6em}}
\newcommand{\exercice}[7]{% numéro, difficulté, titre, origine, notions, identifiant, page du corrigé
  \par\needspace{8\baselineskip}\vspace{1.1em}\phantomsection
  \def\@currentlabel{#1}\label{e:#6}\pdfbookmark[1]{Exercice #1}{e.#6}%
  \noindent{\sffamily\bfseries\large\color{Bleu}Exercice~#1}\quad\niveau{#2}\hfill{\footnotesize\sffamily\color{Gris}#4}\par\nobreak
  \vspace{0.1em}\noindent{\sffamily\bfseries\color{Ink}#3}\par\nobreak
  \vspace{0.1em}\noindent{\footnotesize\sffamily\color{Gris}Notions : #5.\hfill\mbox{corrigé p.~#7}\par}\nobreak
  \vspace{0.35em}}
\newcommand{\renvoi}[1]{\par\nopagebreak\vspace{0.2em}\noindent\hfill{\scriptsize\sffamily\color{Gris}#1}\par}
\newcommand{\place}{\par\vspace{1.6cm plus 0.6cm minus 0.6cm}}
\newcommand{\correction}[4]{% numéro, titre, identifiant, renvoi
  \par\needspace{7\baselineskip}\vspace{1.3em}\phantomsection
  \def\@currentlabel{#1}\label{e:#3}\pdfbookmark[1]{Exercice #1}{c.#3}%
  {\color{Filet}\hrule height 0.4pt}\vspace{0.5em}
  \noindent{\sffamily\bfseries\large\color{Bleu}Exercice~#1}\quad{\sffamily\bfseries\color{Ink}#2}\hfill{\scriptsize\sffamily\color{Gris}#4}\par\nobreak
  \vspace{0.4em}}
\newcommand{\sommaireligne}[3]{\noindent\hbox to 2.2em{\hfill\sffamily\color{Bleu}#1.}\enspace\begin{minipage}[t]{0.8\linewidth}\raggedright#2\end{minipage}\hfill\pageref{e:#3}\par}
\makeatother
\setlength{\parskip}{0.5em plus 0.15em minus 0.1em}

\makeatletter
\newcommand{\probleme}[7]{% numéro, étape, titre, difficulté, origine, durée, identifiant
  \clearpage\phantomsection\label{p:#7}\pdfbookmark[1]{Problème #1 — #2}{p.#7}%
  {\sffamily\small\color{Gris}\textls[150]{PROBLÈME #1\ \ ·\ \ \MakeUppercase{#2}}}\par\nobreak\vspace{0.15em}
  {\sffamily\bfseries\LARGE\color{Ink}#3}\par\nobreak\vspace{0.3em}
  {\color{Ocre}\rule{2.6cm}{1.2pt}}\par\nobreak\vspace{0.2em}
  \noindent{\footnotesize\sffamily\color{Gris}Difficulté~\niveau{#4}\quad·\quad durée indicative : #6\quad·\quad #5}\par\nobreak\vspace{0.3em}}
\newcommand{\fichepb}[3]{% notions, objectifs, prérequis
  {\small\setlength{\parskip}{0.15em}%
  \noindent\hangindent=2.4cm\hangafter=1\makebox[2.4cm][l]{\sffamily\textcolor{Bleu}{Notions}}#1\par
  \noindent\hangindent=2.4cm\hangafter=1\makebox[2.4cm][l]{\sffamily\textcolor{Bleu}{Objectifs}}#2\par
  \noindent\hangindent=2.4cm\hangafter=1\makebox[2.4cm][l]{\sffamily\textcolor{Bleu}{Prérequis}}#3\par}
  \vspace{0.4em}{\color{Filet}\hrule height 0.4pt}\vspace{0.6em}}
\newcommand{\titresection}[2]{% titre, texte
  \clearpage\phantomsection\pdfbookmark[0]{#1}{s.#1}%
  {\sffamily\bfseries\fontsize{24}{28}\selectfont\color{Ink}#1}\par\vspace{0.3em}
  {\color{Ocre}\rule{3.2cm}{1.4pt}}\par\vspace{0.5em}{\small\color{Gris}#2}\par\vspace{0.6em}}
\newcommand{\soustitre}[3]{% numéro, titre, label
  \par\needspace{6\baselineskip}\vspace{1em}\phantomsection\label{#3}%
  {\color{Filet}\hrule height 0.4pt}\vspace{0.5em}
  \noindent{\sffamily\bfseries\large\color{Bleu}Problème~#1}\quad{\sffamily\bfseries\color{Ink}#2}\par\nobreak\vspace{0.3em}}
\makeatother

\begin{document}
\thispagestyle{plain}
{\sffamily\small\color{Gris}\textls[150]{LE COLLÉ · MATHÉMATIQUES PT · CHAPITRE 4}}\par\vspace{0.2em}
{\sffamily\bfseries\fontsize{26}{30}\selectfont\color{Ink}Sélection --- Équations différentielles}\par\vspace{0.3em}
{\sffamily\Large\color{BleuMoyen}Cinq problèmes à travailler après le cours}\par\vspace{0.4em}
{\color{Ocre}\rule{3.2cm}{1.4pt}}\par\vspace{0.8em}
{\small Ces cinq problèmes sont la sélection principale pour les séances, \emph{après} l'explication du cours.
Ils vont de la mise en application au niveau des devoirs et des concours ; chacun réunit plusieurs résultats du chapitre.
Ils complètent le TD complet et le livre, qu'ils ne remplacent pas.

\textbf{Mode d'emploi.} Les traiter dans l'ordre. Pour chaque problème : lire l'en-tête (notions, objectifs, prérequis),
chercher au brouillon, puis rédiger. En cas de blocage, lire \emph{une seule} indication (page des indications), puis reprendre.
Ne lire la correction qu'une fois le problème terminé ou après un vrai effort : elle explique le \emph{choix de la méthode}
avant les calculs. L'origine de chaque problème est indiquée en tête ; les problèmes « d'après » un concours sont des
adaptations rédigées par Le Collé, pas les énoncés officiels.}
\vspace{1.2em}
{\small\renewcommand{\arraystretch}{1.45}\setlength{\tabcolsep}{4pt}\noindent\begin{tabular}{@{}r>{\raggedright\arraybackslash\sffamily\footnotesize}p{2.9cm}>{\raggedright\arraybackslash}p{6.1cm}ll@{\hspace{6pt}}ccc@{}}
& & & & & {\footnotesize\sffamily énoncé} & {\footnotesize\sffamily indic.} & {\footnotesize\sffamily corrigé}\\\hline
1 & \textcolor{Gris}{Mise en application} & Régime transitoire et régime permanent & \niveau{2} & 35\,min & \pageref{p:sel-ed-regimes} & \pageref{i:sel-ed-regimes} & \pageref{c:sel-ed-regimes}\\
2 & \textcolor{Gris}{Méthodes classiques} & Résonance : quand les solutions sont-elles bornées ? & \niveau{3} & 25\,min & \pageref{p:edn-8} & \pageref{i:edn-8} & \pageref{c:edn-8}\\
3 & \textcolor{Gris}{Approfondissement} & Une équation fonctionnelle & \niveau{3} & 20\,min & \pageref{p:edn-18} & \pageref{i:edn-18} & \pageref{c:edn-18}\\
4 & \textcolor{Gris}{Problème de synthèse} & Développement limité d'une solution & \niveau{3} & 30\,min & \pageref{p:edn-7} & \pageref{i:edn-7} & \pageref{c:edn-7}\\
5 & \textcolor{Gris}{Niveau DS / concours} & Solutions périodiques d'un oscillateur forcé & \niveau{5} & 1\,h & \pageref{p:edn-15} & \pageref{i:edn-15} & \pageref{c:edn-15}\\
\end{tabular}}
\probleme{1}{Mise en application}{Régime transitoire et régime permanent}{2}{Problème original Le Collé}{35\,min}{sel-ed-regimes}
\fichepb{premier ordre, équation caractéristique, passage aux complexes, solution périodique, comportement en $+\infty$.}{Résoudre des équations à coefficients constants, puis interpréter la structure « solution particulière + solution homogène » en termes de régimes.}{Résolution du premier ordre ; Théorème sur l'équation caractéristique ; Écriture $a\cos t+b\sin t=A\cos(t-\varphi)$.}
\emph{Partie A : un circuit du premier ordre.} On considère $(E_1)$ : $y'+y=\cos t$.
\begin{enumerate}
\item Résoudre $(E_1)$ sur $\R$.
\item Montrer que $(E_1)$ admet une unique solution $2\pi$-périodique $y_p$, et l'écrire sous la forme $y_p(t)=A\cos(t-\varphi)$ avec $A>0$.
\item Montrer que, pour toute solution $y$ de $(E_1)$, $y(t)-y_p(t)\xrightarrow[t\to+\infty]{}0$. Interpréter : régime transitoire, régime permanent.
\end{enumerate}
\emph{Partie B : un oscillateur amorti.} Soit $\lambda\geq0$ et $(H_\lambda)$ : $y''+2\lambda y'+y=0$.
\begin{enumerate}\setcounter{enumi}{3}
\item Résoudre $(H_\lambda)$ sur $\R$ en distinguant $\lambda=0$, $0<\lambda<1$, $\lambda=1$ et $\lambda>1$.
\item Montrer que, si $\lambda>0$, toutes les solutions de $(H_\lambda)$ tendent vers $0$ en $+\infty$. Que se passe-t-il pour $\lambda=0$ ?
\item On prend $\lambda=\frac12$ et l'on considère $(E_2)$ : $y''+y'+y=\cos t$. Déterminer une solution particulière, puis montrer que toutes les solutions de $(E_2)$ se rapprochent d'une même fonction périodique : leur différence avec elle tend vers $0$ en $+\infty$.
\end{enumerate}
\par\vspace{0.8em}\noindent\hfill{\scriptsize\sffamily\color{Gris}indications p.~\pageref{i:sel-ed-regimes} · correction p.~\pageref{c:sel-ed-regimes}}
\probleme{2}{Méthodes classiques}{Résonance : quand les solutions sont-elles bornées ?}{3}{Exercice complémentaire · TD complet, exercice~14}{25\,min}{edn-8}
\fichepb{équation caractéristique, passage aux complexes, résonance, discussion selon un paramètre.}{Choisir la forme de la solution particulière selon que $i$ est racine ou non, et en tirer le comportement des solutions.}{Équation caractéristique ; Second membre $e^{mt}$ : degré selon la multiplicité de $m$.}
Soit $\omega>0$. Résoudre $y''+\omega^2y=\cos t$, en distinguant $\omega\neq1$ et $\omega=1$. Dans quel cas toutes les solutions
sont-elles bornées ?
\par\vspace{0.8em}\noindent\hfill{\scriptsize\sffamily\color{Gris}indications p.~\pageref{i:edn-8} · correction p.~\pageref{c:edn-8}}
\probleme{3}{Approfondissement}{Une équation fonctionnelle}{3}{Exercice complémentaire · TD complet, exercice~15}{20\,min}{edn-18}
\fichepb{dériver une relation, équation du second ordre, condition nécessaire et réciproque.}{Se ramener à une équation différentielle, puis ne garder que les solutions de l'équation de départ.}{Équation $y''+y=0$ ; Dérivée de $x\mapsto f(-x)$.}
Déterminer les fonctions $f$ dérivables sur $\R$ telles que, pour tout réel $x$, $f'(x)=f(-x)$.
\par\vspace{0.8em}\noindent\hfill{\scriptsize\sffamily\color{Gris}indications p.~\pageref{i:edn-18} · correction p.~\pageref{c:edn-18}}
\probleme{4}{Problème de synthèse}{Développement limité d'une solution}{3}{Exercice complémentaire · TD complet, exercice~16}{30\,min}{edn-7}
\fichepb{premier ordre, développement limité, identification de coefficients.}{Obtenir le développement limité d'une solution sans la calculer, à partir de l'équation.}{Formule de Taylor-Young ; Développement limité de $\sin$.}
Soit $f$ la solution sur $]-\infty,1[$ de $2(t-1)y'+y=\sin(2t)+t^2$ telle que $f(0)=0$. Sans résoudre l'équation, déterminer le
développement limité de $f$ à l'ordre $4$ en $0$.
\par\vspace{0.8em}\noindent\hfill{\scriptsize\sffamily\color{Gris}indications p.~\pageref{i:edn-7} · correction p.~\pageref{c:edn-7}}
\probleme{5}{Niveau DS / concours}{Solutions périodiques d'un oscillateur forcé}{5}{Exercice complémentaire · TD complet, exercice~17}{1\,h}{edn-15}
\fichepb{intégrale fonction de sa borne, problème de Cauchy, périodicité, résonance.}{Exprimer une solution par une intégrale, puis caractériser l'existence de solutions périodiques, comme dans un problème de concours.}{Dérivée de $x\mapsto\int_0^xg$ ; Intégrale d'une fonction périodique sur une période.}
Soit $f$ une fonction continue sur $\R$ et l'équation $(E)$ : $y''+y=f(t)$.
\begin{enumerate}
\item Montrer que $y_0(t)=\displaystyle\int_0^t\sin(t-u)f(u)\,du$ est la solution de $(E)$ telle que $y_0(0)=y_0'(0)=0$.
\emph{(Développer $\sin(t-u)$ pour dériver.)}
\item On suppose $f$ $2\pi$-périodique, et l'on pose $C=\int_0^{2\pi}f(u)\cos u\,du$, $S=\int_0^{2\pi}f(u)\sin u\,du$. Montrer que
$y_0(t+2\pi)-y_0(t)=C\sin t-S\cos t$.
\item En déduire que $(E)$ admet une solution $2\pi$-périodique si et seulement si $C=S=0$, et qu'alors toutes les solutions
le sont.
\item Application : $f(t)=\cos t$, puis $f(t)=\cos(2t)$.
\end{enumerate}
\par\vspace{0.8em}\noindent\hfill{\scriptsize\sffamily\color{Gris}indications p.~\pageref{i:edn-15} · correction p.~\pageref{c:edn-15}}
\titresection{Indications}{Une indication par question délicate. N\'en lire qu\'une à la fois, puis reprendre la recherche.}
\soustitre{1}{Régime transitoire et régime permanent}{i:sel-ed-regimes}
\begin{itemize}[leftmargin=1.2em,itemsep=0.25em]\item 1. Résoudre $z'+z=e^{it}$, puis prendre la partie réelle.
\item 2. Si $K\neq0$, la fonction $Ke^{-t}$ est strictement monotone : peut-elle être périodique ?
\item 4. Discuter le signe de $\lambda^2-1$ ; pour $\lambda>1$, regarder la somme et le produit des racines.
\item 5. Toutes les exponentielles ont une partie réelle strictement négative.
\item 6. $i^2+i+1=i$ : $z=\frac{e^{it}}i$.
\end{itemize}
\soustitre{2}{Résonance : quand les solutions sont-elles bornées ?}{i:edn-8}
\begin{itemize}[leftmargin=1.2em,itemsep=0.25em]\item Résoudre $z''+\omega^2z=e^{it}$ : $i$ est-il racine de $r^2+\omega^2=0$ ?
\item Si $\omega=1$ : chercher $z=Q(t)e^{it}$ avec $\deg Q=1$ ; $Q''+2iQ'=1$.
\item Second membre $\cos t$ : prendre la partie \emph{réelle}.
\end{itemize}
\soustitre{3}{Une équation fonctionnelle}{i:edn-18}
\begin{itemize}[leftmargin=1.2em,itemsep=0.25em]\item Montrer que $f$ est deux fois dérivable, puis dériver la relation.
\item $f''=-f$ n'est qu'une condition nécessaire : reporter $A\cos+B\sin$ dans la relation de départ.
\end{itemize}
\soustitre{4}{Développement limité d'une solution}{i:edn-7}
\begin{itemize}[leftmargin=1.2em,itemsep=0.25em]\item Justifier que $f$ est de classe $\mathcal C^\infty$ : $f'$ s'exprime à l'aide de $f$.
\item Écrire $f=\sum a_kt^k+o(t^5)$ et identifier le coefficient de $t^k$ dans $2(t-1)f'+f$.
\item $a_0=f(0)=0$ ; on trouve successivement $a_1$, $a_2$, $a_3$, $a_4$.
\end{itemize}
\soustitre{5}{Solutions périodiques d'un oscillateur forcé}{i:edn-15}
\begin{itemize}[leftmargin=1.2em,itemsep=0.25em]\item 1. Développer $\sin(t-u)=\sin t\cos u-\cos t\sin u$ et sortir de l'intégrale ce qui ne dépend pas de $u$.
\item 2. $\int_t^{t+2\pi}g=\int_0^{2\pi}g$ pour $g$ continue et $2\pi$-périodique.
\item 3. Une fonction $C\sin t-S\cos t$ est nulle si et seulement si $C=S=0$ : l'évaluer en $0$ et en $\frac\pi2$.
\item 4. $\int_0^{2\pi}\cos^2u\,du=\pi$ ; $\int_0^{2\pi}\cos2u\cos u\,du=0$.
\end{itemize}
\titresection{Corrections}{Chaque correction commence par le choix de la méthode ; les encadrés résument ce qui resservira. Les renvois « méthode~1.4 du livret » désignent le livret du chapitre.}
\soustitre{1}{Régime transitoire et régime permanent}{c:sel-ed-regimes}
\choix{au premier comme au second ordre : solutions de l'équation homogène, puis une solution particulière de la forme du second membre (passage aux complexes). Le comportement en $+\infty$ se lit sur les exponentielles de l'équation homogène.}
\q{1.} Homogène : $Ke^{-t}$. On résout $z'+z=e^{it}$ : $z=ce^{it}$ avec $(i+1)c=1$, $c=\frac{1-i}2$, et
$\operatorname{Re}\left(\frac{1-i}2e^{it}\right)=\frac{\cos t+\sin t}2$. Solutions : $y=\frac{\cos t+\sin t}2+Ke^{-t}$.

\q{2.} $y_p=\frac{\cos t+\sin t}2$ est $2\pi$-périodique. Si $K\neq0$, $y-y_p=Ke^{-t}$ n'est pas périodique (elle est strictement monotone), donc $y$ ne l'est pas : $y_p$ est l'unique solution périodique. Avec $\cos t+\sin t=\sqrt2\cos\left(t-\frac\pi4\right)$ : $y_p=\frac{\sqrt2}2\cos\left(t-\frac\pi4\right)$, $A=\frac{\sqrt2}2$, $\varphi=\frac\pi4$.

\q{3.} $y-y_p=Ke^{-t}\to0$. Le terme $Ke^{-t}$, qui dépend de la condition initiale, s'efface : c'est le \emph{régime transitoire} ; il reste la réponse périodique $y_p$, indépendante de la condition initiale : le \emph{régime permanent}.

\q{4.} Équation caractéristique $r^2+2\lambda r+1=0$, de discriminant réduit $\lambda^2-1$.
\begin{itemize}
\item $\lambda=0$ : $r=\pm i$, $y=A\cos t+B\sin t$.
\item $0<\lambda<1$ : $r=-\lambda\pm i\omega$ avec $\omega=\sqrt{1-\lambda^2}$, $y=e^{-\lambda t}(A\cos\omega t+B\sin\omega t)$ (régime pseudo-périodique).
\item $\lambda=1$ : racine double $-1$, $y=(A+Bt)e^{-t}$ (régime critique).
\item $\lambda>1$ : deux racines réelles $r_\pm=-\lambda\pm\sqrt{\lambda^2-1}$, toutes deux strictement négatives (leur produit vaut $1$ et leur somme $-2\lambda<0$), $y=Ae^{r_+t}+Be^{r_-t}$ (régime apériodique).
\end{itemize}

\q{5.} Si $\lambda>0$, dans chaque cas les solutions sont combinaisons de fonctions $t^ke^{\alpha t}\cos(\beta t)$ ou $t^ke^{\alpha t}\sin(\beta t)$ avec $\alpha<0$ : elles tendent vers $0$ (croissances comparées pour $\lambda=1$). Pour $\lambda=0$, les solutions non nulles sont périodiques et ne tendent pas vers $0$ : pas d'amortissement.

\q{6.} On résout $z''+z'+z=e^{it}$ : $i$ n'est pas racine ($i^2+i+1=i\neq0$), $z=\frac{e^{it}}i=-ie^{it}$, de partie réelle $\sin t$. Donc $y_p=\sin t$ et, d'après 4 et 5 ($\lambda=\frac12$), toute solution s'écrit $y=\sin t+h(t)$ avec $h$ solution de $(H_{1/2})$, qui tend vers $0$ : $y(t)-\sin t\to0$.
\retenir{équation linéaire à coefficients constants, amortie, forcée par un signal périodique : la solution générale de l'équation homogène est le régime transitoire, la solution particulière périodique est le régime permanent.}
\par\hfill\textcolor{BleuClair}{$\square$}\par
\soustitre{2}{Résonance : quand les solutions sont-elles bornées ?}{c:edn-8}
$S_H=\{A\cos\omega t+B\sin\omega t\}$. On résout $z''+\omega^2z=e^{it}$.
\begin{itemize}
\item $\omega\neq1$ : $i$ n'est pas racine de $r^2+\omega^2$, $z=\frac{e^{it}}{\omega^2-1}$, d'où $y_p=\frac{\cos t}{\omega^2-1}$.
Toutes les solutions sont bornées.
\item $\omega=1$ : $i$ est racine simple, $z=Q(t)e^{it}$ avec $Q''+2iQ'=1$, $Q=-\frac i2t$ ; $y_p=\operatorname{Re}z=\frac t2\sin t$.
Aucune solution n'est bornée (le terme $\frac t2\sin t$ ne l'est pas).
\end{itemize}
\retenir{excitation à la pulsation propre : résonance, amplitude non bornée.}
\par\hfill\textcolor{BleuClair}{$\square$}\par
\soustitre{3}{Une équation fonctionnelle}{c:edn-18}
\choix{on se ramène à une équation différentielle en dérivant la relation.}
Soit $f$ une solution. Comme $x\mapsto f(-x)$ est dérivable, $f'$ l'est : $f$ est deux fois dérivable et
$f''(x)=-f'(-x)=-f(x)$. Donc $f=A\cos+B\sin$. Réciproquement, $f'(x)=-A\sin x+B\cos x$ et $f(-x)=A\cos x-B\sin x$ :
l'égalité pour tout $x$ équivaut à $B=A$. Les solutions sont les fonctions $x\mapsto A(\cos x+\sin x)$.
\piege{l'équation $f''=-f$ n'est qu'une condition \emph{nécessaire} : la réciproque élimine une constante.}
\par\hfill\textcolor{BleuClair}{$\square$}\par
\soustitre{4}{Développement limité d'une solution}{c:edn-7}
\choix{on n'a pas besoin de $f$ : l'équation fournit une relation entre les coefficients du développement.}
Sur $]-\infty,1[$, $f'=\frac{\sin2t+t^2-f}{2(t-1)}$ ; par récurrence, $f$ est de classe $\mathcal C^\infty$, donc admet un
développement $f(t)=\sum_{k=0}^5a_kt^k+o(t^5)$, et $f'(t)=\sum_{k=1}^5ka_kt^{k-1}+o(t^4)$. Le coefficient de $t^k$ dans
$2tf'-2f'+f$ est $(2k+1)a_k-2(k+1)a_{k+1}$, et $\sin2t+t^2=2t+t^2-\frac43t^3+o(t^4)$. Avec $a_0=f(0)=0$ :
\[a_0-2a_1=0,\quad3a_1-4a_2=2,\quad5a_2-6a_3=1,\quad7a_3-8a_4=-\tfrac43,\]
d'où $a_1=0$, $a_2=-\frac12$, $a_3=-\frac7{12}$, $a_4=-\frac{11}{32}$ :
\[f(t)=-\frac{t^2}2-\frac{7t^3}{12}-\frac{11t^4}{32}+o(t^4).\]
\par\hfill\textcolor{BleuClair}{$\square$}\par
\soustitre{5}{Solutions périodiques d'un oscillateur forcé}{c:edn-15}
\q{1.} $y_0(t)=\sin t\int_0^tf(u)\cos u\,du-\cos t\int_0^tf(u)\sin u\,du$. En dérivant (produit et intégrale de borne variable),
les termes $\sin t\,f(t)\cos t-\cos t\,f(t)\sin t$ s'annulent : $y_0'(t)=\cos t\int_0^tf\cos+\sin t\int_0^tf\sin$. En dérivant
encore : $y_0''=-y_0+f(t)(\cos^2t+\sin^2t)=-y_0+f$. Enfin $y_0(0)=y_0'(0)=0$.

\q{2.} Comme $f(u)\cos u$ est $2\pi$-périodique, $\int_0^{t+2\pi}f\cos=\int_0^tf\cos+\int_t^{t+2\pi}f\cos=\int_0^tf\cos+C$, et de même avec
$\sin$ ; en reportant dans l'expression de $y_0$, on obtient $y_0(t+2\pi)=y_0(t)+C\sin t-S\cos t$.

\q{3.} Les solutions sont $y=y_0+A\cos t+B\sin t$, et $y(t+2\pi)-y(t)=C\sin t-S\cos t$ pour toutes. Cette fonction est nulle si et
seulement si $C=S=0$ (prendre $t=0$ et $t=\frac\pi2$). Donc ou bien toutes les solutions sont $2\pi$-périodiques, ou bien aucune.

\q{4.} $f=\cos$ : $C=\pi\neq0$, aucune solution périodique (résonance : on retrouve $\frac t2\sin t$). $f(t)=\cos2t$ : $C=S=0$, toutes
les solutions sont $2\pi$-périodiques (par exemple $y_p=-\frac13\cos2t$).
\par\hfill\textcolor{BleuClair}{$\square$}\par
\end{document}
